\documentclass[svgnames]{article}     % use "amsart" instead of "article" for AMSLaTeX format
%\geometry{landscape}                 % Activate for rotated page geometry

%\usepackage[parfill]{parskip}        % Activate to begin paragraphs with an empty line rather than an indent

\usepackage{graphicx}                 % Use pdf, png, jpg, or eps§ with pdflatex; use eps in DVI mode

%maths                                % TeX will automatically convert eps --> pdf in pdflatex
\usepackage{amssymb}
\usepackage{amsmath}
\usepackage{esint}
\usepackage[bottom = 0.3in, top = 0.3in]{geometry}

% Inverting Color of PDF
%\usepackage{xcolor}
%\pagecolor[rgb]{0.19,0.19,0.19}
%\color[rgb]{0.77,0.77,0.77}

%noindent
%\setlength\parindent{0pt}

%pgfplots
\usepackage{pgfplots}
\usepackage{ amssymb }

%images
\graphicspath{{/Users/devaldeliwala/research/AQT/review/}}                   % Activate to set a image directory

%tikz
\usepackage{pgfplots}
\newcommand{\Tau}{\mathrm{T}}
\pgfplotsset{compat=1.15}
\usepackage{comment}
\usetikzlibrary{arrows}
\usepackage[most]{tcolorbox}

%Figures
\usepackage{float}
\usepackage{caption}
\usepackage{lipsum}


\title{Introduction to Superconducting Qubit Engineering}
\author{Deval Deliwala}
%\date{}                              % Activate to display a given date or no date

\begin{document}
\maketitle
%\section{}
%\subsection{}
%\tableofcontents                     % Activate to display a table of contents

In superconducting qubit processors, information is stored in the quantum
degrees of freedom of nanofabricated, \textit{anharmonic} oscillators
constructed from superconducting circuit elements. The energy-level spectra of
s.c. qubits are governed by circuit element parameters and thus are
configurable. They have a rich parameter space of possible qubit properties and
operation regimes, with a predictable performance, tunable to exhibit
``atom-like'' properties -- an \textit{artificial atom}. 

\section{Engineering Quantum Circuits}

In this section, we will demonstrate how s.c. quantum systems can be engineered
to achieve certain desired properties.  

\subsection{Quantum Harmonic Oscillator $\rightarrow$ Transmon Qubit}

A quantum system is governed by the time-dependent Schr\"odinger equation, 

\[
  \hat{H}|\psi(t)\rangle = i\hbar \frac{\partial }{\partial t} |\psi(t)\rangle,
\] \vspace{3px}

where $|\psi(t)\rangle$ is the state of the quantum system at time $t$ and
$\hat{H}$ is the Hamiltonian that describes the total energy of the system.
Solving the differential equation yields 

\[
  |\psi(t)\rangle = e^{-i\hat{H} t / \hbar} |\psi(0)\rangle. 
\] \vspace{3px}

The time-independent Hamiltonian $\hat{H}$ is the generator for time-evolution
via the operator $e^{-i \hat{H} t / \hbar}$. Determining the Hamiltonian of
a system is therefore the first step to deriving its dynamical behavior. To understand the dynamics of a s.c. qubit circuit, we start with the classical
description of a linear LC resonant circuit. 

\begin{figure}[H]
  \centering
    \includegraphics[width = 4cm]{screenshot.png}
    \caption{Simple LC resonant circuit.}
\end{figure}


Here, the energy oscillates between electrical energy in the capacitor $C$ and
magnetic energy in the inductor $L$. We will arbitrarily associate the
electrical energy with the ``kinetic energy" and the magnetic energy with the
``potential energy" of the oscillator. The instantaneous, time-dependent energy
in each element is derived from its current and voltage, 

\begin{align} \label{energy} 
E(t) = \int_{-\infty}^{t} V(t')I(t') \, dt', 
\end{align} \vspace{3px}

where $V(t')$ and $I(t')$ denote the voltage and current of the capacitor or
inductor. To derive the classical Hamiltonian, we use the Lagrange-Hamilton
formulation. Here, we represent the circuit elements in terms of one of its
generalized circuit coordinates -- charge or flux. We choose flux, defined as
the time integral of the voltage 

\[
\Phi(t) = \int_{-\infty}^{t} V(t') \, dt'. 
\] \vspace{3px}

Here, the voltage at the node is also the branch voltage across the element.
Combining the above equations for $\Phi(t)$ and $E(t)$, and using the relations 

\[
V = L \frac{dI}{dt} \qquad I = C \frac{dV}{dt}, 
\] \vspace{3px}

we write down energy terms for the capacitor and inductor in terms of the node
flux,

\begin{align} \label{}
  \tau_C &= \frac{1}{2} C \dot{\Phi}^2, \\ 
  U_L &= \frac{1}{2L} \Phi^2. 
\end{align}\vspace{3px}

The Lagrangian is defined as the difference between the kinetic and potential
energy terms and can thus be expressed as 

\[
  \mathcal{L} = \tau_C - U_L = \frac{1}{2}C \dot{\Phi}^2 - \frac{1}{2L}\Phi^2. 
\] \vspace{3px}

Then, using the Legendre transformation, for which we need to calculate the
momentum conjugate to the flux, which in this case is the charge on the
capacitor, 

\[
  Q = \frac{\partial \mathcal{L}}{\partial \dot{\Phi}} = C \dot{\Phi}
\] \vspace{3px}

The Hamiltonian of the system is now defined as 

\begin{align}
  H = Q\dot{\Phi} - \mathcal{L} = \frac{Q^2}{2C} + \frac{\Phi^2}{2L} \equiv
  \frac{1}{2}CV^2 + \frac{1}{2}LI^2, 
\end{align} \vspace{3px}


as one would expect for an electrical LC circuit. This Hamiltonian is
analogous  to that of a mechanical harmonic oscillator with mass $m = C$ and
a resonant frequency $\omega = 1 / \sqrt{LC}$, which expressed in position $x$,
and momentum, $p$ coordinates takes the form $H = \frac{p^2}{2m} + m\omega^2
x^2 / 2$.

However, the Hamiltonian described above is classical. To proceed to
a quantum-mechanical description, we promote the charge and flux coordinates to
quantum operators. And whereas the classical coordinates satisfy the Poisson
bracket: 

\begin{align*}
  \{f, g\} &= \frac{\delta f}{\delta \Phi} \frac{\delta g}{\delta Q}
  - \frac{\delta g}{\delta \Phi} \frac{\delta f}{\delta Q} \\ 
  \rightarrow \{ \Phi, Q \} &= \frac{\delta \Phi}{\delta \Phi} \frac{\delta
  Q}{\delta Q} - \frac{\delta Q}{\delta \Phi} \frac{\delta \Phi}{\delta Q}
  = 1 - 0 = 1, 
\end{align*}

the quantum operators similarly satisfy a commutation relation: 

\begin{align} \label{}
  [ \hat{\Phi}, \hat{Q} ] = \hat{\Phi} \hat{Q} - \hat{Q} \hat{\Phi} = i\hbar
\end{align}\vspace{3px}

Notably the same as the commutation of position and momentum, $[\hat{x},
\hat{p}]$. We can take $\hat{\Phi} \rightarrow \Phi$ and $\hat{Q} \rightarrow
i\hbar \frac{d }{d \Phi} $ to yield the quantum operators $\hat{Q}$ and $
\hat{\Phi}$ analogous to position and momentum. 

In a simple LC resonant circuit, shown in Figure 1, both the inductor $L$ and
the capacitor $C$ are linear circuit elements. Defining the reduced flux $\phi
\equiv 2\pi \Phi / \Phi_0$ and the reduced charge $n = Q / 2e$, we can write
down the following quantum-mechanical Hamiltonian for the circuit, 

\begin{align}\label{QHO}
  H = 4E_C n^2 + \frac{1}{2}E_L \phi^2
\end{align}

where $E_C = e^2 / (2C)$ is the charging energy required to add \textit{each}
electron of the Cooper-pair to the island and $E_L = (\Phi_0 / 2\pi)^2 / L$ is
the inductive energy, where $\Phi_0 = h / (2e)$ is the superconducting magnetic
flux quantum. The quantum operator  $n$ is the excess number of Cooper-pairs on
the island, and $\phi$ -- the reduced flux -- is denoted the ``gauge-invariant
phase'' across the inductor. These two operators obey the commutation relation
$[\phi, n] = i$. The factor 4 in front of the charging energy $E_c$ is purely
a historical artifact, first defined for single-electron systems and afterwards
adopted to two-electron Cooper-pair systems. 

The Hamiltonian in Equation \ref{QHO}. is identical to the quantum harmonic oscillator.
We can treat $\phi$ as the generalized position coordinate making the first
term the kinetic energy and the second term potential energy. The functional form of the potential energy influences the possible energy
eigenstates. The fact that the potential energy term is quadratic ($U_L \propto
\phi^2$) in Equation \ref{QHO}. gives rise to the shape of the potential below. 

\begin{figure}[H]
  \centering
    \includegraphics[width = 5cm]{screenshot 1.png}
    \caption{Energy potential for the quantum harmonic oscillator, whose energy
    eigenstates are equidistantly spaced $\hbar \omega_r$ apart.} 
\end{figure}

The solution to the eigenvalue problem gives an infinite series of possible
eigenstates $|k\rangle, (k = 0, 1, 2, \hdots)$, whose corresponding
eigenenergies $E_k$ are all equidistantly spaced, i.e. $E_{k+1} - E_k
= \hbar\omega_r$, where $\omega_r = \sqrt{8E_LE_C} / \hbar = 1/\sqrt{LC}$
denotes the resonant frequency of the system. We can represent these results in
a more compact form for the quantum harmonic oscillator (QHO) Hamiltonian using
creation and annihilation operators: 

\begin{align}\label{creationQHO}
H = \hbar \omega_r \left( a^\dagger a + \frac{1}{2} \right)  
\end{align} \vspace{3px}

where $a^\dagger a$ is the creation (annihilation) operator of a single
excitation of the resonator. The Hamiltonian in Equation \ref{creationQHO}. is written as an
energy. From here on, $\hbar$ will be omitted; using natural units. 

The linear characteristics of the QHO has a liimitation in its applications for
quantum information processing. We need to be able to define a computational
subspace consisting of only two energy states (usually the two lowest energy
eigenstates) in between which transitions can be driven without also exciting
other levels in the system. Because many gate operations depend on frequency
selectivity, the equidistant level-spacing of the QHO is not possible to work
with. 

Therefore, we add anharmonicity (nonlinearity) into our system. We require the
transition frequencies $\omega_q^{0\rightarrow 1}$ and $\omega_q^{1\rightarrow
2} $ to be different enough to be individually addressable. In general, the
   larger the anharmonicity the better, however, the amount affects how short
   pulses used to drive the qubit can be. 

To introduce anharmonicity, we replace the linear inductor of the QHO with
a Josephson junction -- a nonlinear, dissipationless circuit element that forms
the backbone of superconducting circuits. The potential energy of the Josephson
junction can be derived similarly using Equation \ref{energy}. The two
Josephson junction relations 

\[
I = I_c \sin(\phi), \quad V = \frac{\hbar}{2e} \frac{d \phi}{d t},
\] \vspace{3px}

result in the modified Hamiltonian 

\begin{align} \label{joseph}
  H = 4E_C n^2 - E_J \cos(\phi)
\end{align}\vspace{3px}

where $E_C = e^2 / (2C_\Sigma), C_\Sigma = C_s + C_J$ is the total capacitance,
including both shunt capacitance $C_s$ and the self-capacitance of the junction
$C_J$, and $E_J = I_C\Phi_0 / 2\pi$ is the Josephson energy, with $I_C$ being
the critical current of the junction. 

The potential energy no longer takes a pure parabolic form, but rather
features a cosinusoidal form from the second term in Equation \ref{joseph}.,
which makes the energy spectrum nondegenerate. This allows us to make the
oscillator anharmonic and thus allows us to identify a uniquely addressable
two-level quantum system shown below. 

\begin{figure}[H]
  \centering
    \includegraphics[width = 5cm]{screenshot 2.png}
    \caption{The Josephson junction reshapes the quadratic energy potential
      (dashed red) into sinusoidal (solid blue), which yields non-equidistant
      energy levels. This allows us to isolate the two lowest energy levels
      $|0\rangle$ and $| 1 \rangle$, forming a computational basis with energy
      separation $\hbar \omega_{01}$, which is different than $\hbar
    \omega_{12}$.} 
\end{figure}

The system dynamics are governed by the dominant energy in Equation
\ref{joseph}., reflected in the energy $E_J / E_C$ ratio. The s.c. qubit
community has converged towards circuitry with $E_J \gg E_C$, as  $E_J \leq
E_C$ results in qubits becoming highly sensitive to charge noise which is more
difficult to deal with than flux noise. Additionally there is more flexibility
in engineering the inductive (potential) part of the Hamiltonian.

To make $E_J \gg E_C$, one proffered approach is to make $E_C$ small by
shunting the junction with a large capacitor, $C_s \gg C_J$, effectively making
the qubit less sensitive to charge noise -- a circuit commonly known as the
\textit{transmon} qubit. Here, the superconducting phase $\phi$ can be used as
a quantum number because the spread (or \textit{quantum fluctuation}) of $\phi$
values represented by the wave function is small. The low-energy eigenstates
are therefore, to a good approximation, localized states in the potential well. 

If excitation to higher non-computational states is suppressed over any gate
operations, either due to a large enough $|\alpha| = |\omega_q^{1\rightarrow 2}
- \omega_q^{0\rightarrow 1}|$  or due to any robust control techniques, we can
effectively treat the anharmonic oscillator (AHO) as a quantum two-level
system, simplifying the Hamiltonian to 

\begin{align} \label{twoh}
  H = \omega_q \frac{\sigma_z}{2} = \omega_q \begin{pmatrix}
    \frac{1}{2} & 0 \\ 0 & -\frac{1}{2}
  \end{pmatrix} 
\end{align}\vspace{3px}

where $\sigma_z$ is the Pauli-z operator. However, higher levels are often
useful to implement more efficient gate operations and should not be ignored.

\subsection{Qubit Hamiltonian Engineering} 

\subsubsection{Tunable qubit: split transmon} 

To implement fast gate operations with high-fidelity, quantum processor
architectures feature tunable qubit frequencies. For instance, in some cases,
we need to bring two qubits into resonance to exchange (SWAP Gate) energies,
while also having the capability of separating them during idling periods to
minimize their interactions. To do this, we need an external parameter which
allows us to access one of the degrees of freedom of the system in
a controllable fashion.  

One common technique is to replace the single Josephson junction with a loop
interrupted by two identical junctions -- forming a DC superconducting quanutm
interference device (dc-SQUID). Due to the interference between the two arms of
the SQUID, the effective critical current of the two parallel junctions can be
decreased by applying a magnetic flux threading the loop. 

\begin{figure}[H]
  \centering
    \includegraphics[width = 4cm]{screenshot 3.png}
  \caption{SQUID modular qubit circuit representation for capacitively shunted
  qubit modalities. Symmetric transmon qubit with Josephson energy $E_J$ is
shunted with a capacitor yielding a charging energy $E_C$.}
\end{figure}

\begin{figure}[H]
  \centering
    \includegraphics[width = 4cm]{screenshot 4.png}
    \caption{Corresponding qubit transition frequencies for the two lowest
      energy states as a function of applied magnetic flux in units of
    $\Phi_0$.} 
    \label{tunable} 
\end{figure}

The algebraic sum of branch flux of all the inductive elements along the loop
plus the externally applied flux equals an integer number of superconducting
flux quanta to match the fluxoid quantization condition. That is, 

\[
\varphi_1 - \varphi_2 + 2\varphi_e = 2\pi k 
\] \vspace{3px}

where $\varphi_e = \pi \Phi_\text{ext} / \Phi_0$. Using this condition, we can
eliminate one degree of freedom and treat the SQUID-loop as a single junction,
but with the important modification that $E_J$ is tunable (via the SQUID
critical current) by means of the external flux $\Phi_\text{ext}$. The
effective Hamiltonian of the so-called ``split transmon" is then 


\begin{align} \label{splittransmon}
  H &= 4E_C n^2 - 2E_J |\cos(\varphi_e)| \cos (\phi) \\ 
  H &= 4E_C n^2 - E_J'(\varphi_e) \cos(\phi)
\end{align}\vspace{3px}

This equation is analogous to Equation \ref{joseph}. with $E_J$ replaced by
$E_J'(\varphi_e) = 2 E_J |\cos (\varphi_e)|$. The magnitude of the net,
effective Josephson energy $E_J'$ has a period of $\Phi_0$ in applied flux and
spans from 0 to its maximum value $2E_J$. Therefore, the qubit frequency can be
tuned periodically with $\Phi_\text{ext} $, shown in Figure \ref{tunable}. 

However, the split transmon also introduces sensitivity to random flux
fluctuations, known as \textit{flux noise}. At any point, the slope of the
qubit spectrum, $\partial \omega_q / \partial \Phi_\text{ext}$ indicates to first
order how strongly this noise affects the qubit frequency. This sensitivity is
only 0 at multiples of the flux quantum, $\Phi_\text{ext} = k\Phi_0$ where $k$
is an integer, $\partial \omega_q / \partial \Phi_\text{ext}  = 0$. 


One method of reducing flux noise sensitivity while maintaining sufficient
tunability is to make the two junctions in the split transmon
\textit{asymmetric}, which yields the following Hamiltonian 

\begin{align} \label{asymH}
  H = 4E_C n^2 - E_{J\Sigma} \sqrt{\cos^2(\varphi_e) + d^2\sin^2(\varphi_e)}
  \cos(\phi)    
\end{align}\vspace{3px}

where $E_{J\Sigma} = E_{ J1} + E_{J2}$ and $d = (\gamma - 1) / (\gamma + 1)$ is
the junction asymmetry parameter, with $\gamma = E_{J2} / E_{J1}$. Going from
symmetric to asymmetric transmons does not change the circuit topology, but has
a profound impact for practical applications. As you can see below, the flux
sensitivity is suppressed across the entire tunable frequency range for
asymmetric circuits relative to symmetric ones. 

\begin{figure}[H]
  \centering
    \includegraphics[width = 10cm]{screenshot 5.png}
    \caption{Symmetric transmon qubit from Figure \ref{tunable} compared to an
      Asymmetric transmon qubit with junction asymmetry $\gamma = E_{J2}
    / E_{J1} = 2.5$.}
\end{figure}


\subsubsection{Flux Qubit and Fluxonium} 


Split transmon qubits, whether symmetric or antisymmetric, still have the same
topology yielding a sinusoidal potential. The degree to which properties of
these qubits can be engineered has not changed. The limited anharmonicity in
transmon-type qubits intrinsically causes significant residual excitation to
higher-energy states, undermining fidelity of gate operations. It's necessary
to introduce more complexity in the circuit. 

The \textit{flux qubit} is engineered by making the qubit loop interrupted by
three or four junctions. 

\begin{figure}[H]
  \centering
    \includegraphics[width = 4cm]{screenshot 6.png}
    \caption{Capacitively shunted flux qubit, where a small principle junction
    (red) is shunted with a capacitor yielding a charging energy $E_C$.}
\end{figure}

On one branch is a smaller junction, and on the other branch are two identical
junctions, both a factor $\gamma$ larger in size.  The addition of one more
junction is non-trivial as it changes the circuit topology and reshapes the
potential energy profile. 


Each junction is associated with a \textit{phase variable} and the fluxoid
quantization condition again allows us to eliminate one degree of freedom. Now,
we have a two-dimensional potential landscape, which compared to the simpler
topology of the transmon, complicates the problem both conceptually and
computationally. Under a \textit{quasi-1D approximation}, the Hamiltonian is 

\begin{align} \label{1dapproxH}
  H \approx 4E_C n^2 - E_J \cos (2\pi + \varphi_e) - 2\gamma E_J
  \cos(\varphi_e)
\end{align}

Note the phase variable in Equation \ref{1dapproxH}. is the sum of the branch
phases across the two array junctions, $\phi = (\varphi_1 + \varphi_2)/ 2$,
assuming same current direction across $\varphi_1$ and $\varphi_2$. The final
potential profile as well as the corresponding energy eigenstates depends on
both the external flux $\varphi_e$ and the junction area ratio $\gamma$. 

Most commonly, the optimal working point is when $\varphi_e = \pi + 2\pi k$,
where $k$ is an integer -- that is when half a superconducting flux quantum
threads the loop. Here, the qubit spectrum reaches its minimum, referred to as
the \textit{flux degeneracy point}, where flux qubits tend to have the optimal
coherence time. Here, the potential energy may assume a single-well ($\gamma
\geq 2$ ) or a double well $(\gamma < 2)$ profile. 

Compared to the flux qubit, which usually contains two or three array
junctions, the number of array junctions in the \textit{fluxonium} qubit
dramatically increases. With the same quasi-1D approximation, the last term in
Equation \ref{1dapproxH} becomes $-N_\gamma E_J \cos (\phi / N)$, giving 

 \begin{align} \label{fluxonium}
  H \approx 4E_C n^2 - E_J \cos (2\phi + \varphi_e) - N\gamma E_J \cos(\phi
  / N)
\end{align}\vspace{3px}

where $N$ denotes the number of array junctions. For large $N$, the argument in
the cosine term,  $\phi / N$ becomes sufficiently small that a second order
expansion is a good approximation. This results in the final fluxonium
Hamiltonian, 

\begin{align} \label{fluxoniumf}
  H \approx 4E_C n^2 - E_J \cos (\phi + \varphi_e) + \frac{1}{2}E_L \phi^2
\end{align}\vspace{3px}

where  $E_L = (\gamma / N)E_J$ is the inductive energy of the
\textit{effective} inductance contributed by the junction array -- often known
as superinductance. The potential energy is thus a quadratic term modulated by
a sinusoidal term, similar to that of an rf-SQUID type flux qubit. However, the
kinetic inductance of the Josephson Junction array is in general much larger
than the geometric inductance of the wire in an rf-SQUID. 

Depending on the relative magnitude of $E_J$ and $E_L$, the fluxonium system
could involve plasmon states (in the same well) and fluxon states (in different
wells). 


\end{document}

